% Chapter 5: RESULTS
% Six sections, each anchored to a single named data source.
% Data freeze: 2026-05-09. All numbers below are a snapshot at that date.

\section{Evaluation Framework}
\label{sec:evaluation_framework}

This chapter reports observed values from the deployed Clothie system. The framework, validity gates, and cohort rules summarised below were defined in Chapter~3, Section~3.6; they are retained here so that every empirical block in §5.2--§5.6 inherits the same boundaries.

\textbf{Data freeze.} All aggregates in §5.2--§5.6 are computed from a snapshot of the production PostgreSQL database taken at \textbf{2026-05-09}. The cohort window is open, that is, no $t_{start}$ filter is applied; sessions from the earliest deployed traffic on 2026-04-20 up to and including 2026-05-08 are eligible. Because the deployed system continues to receive traffic, any reproduction with a later freeze date will return a larger cohort whose distribution may differ; the numbers below are therefore a point-in-time observation, not a stable historical record.

\textbf{Reproducibility note.} Each subsection in §5.2--§5.6 names the PostgreSQL table or view from which its numbers were drawn. A reviewer can reproduce any cell of any table by running the corresponding aggregation against the live database with a freeze-date filter; cell values may shift as the deployed system continues to log new sessions.

\subsection{Inclusion, Exclusion, and Time-Window Rules}
\label{subsec:inclusion_exclusion_rules}

To ensure reproducibility under the open-window cohort, every aggregate in this chapter follows the same boundary policy:
\begin{enumerate}
    \item Include only rows whose \texttt{created\_at} (or equivalent timestamp column) is no later than the data-freeze date 2026-05-09.
    \item Exclude sessions or events with missing key fields for the target metric, for example sessions with no \texttt{llm\_token\_usage} row are excluded from token aggregates but kept in the raw cohort count.
    \item Exclude duplicate interaction rows when the same session-event key appears multiple times due to retries (the unique constraints on \texttt{selected\_items} and \texttt{product\_intents} make this exclusion a no-op in the present cohort).
    \item Report both the raw cohort size and the post-filter cohort size whenever a metric is computed on a subset, so the impact of exclusion is explicit.
\end{enumerate}

\subsection{Data Quality and Validity Checks}
\label{subsec:data_quality_validity}

\begin{table}[H]
\centering
\begin{tabular}{p{4.0cm}p{4.7cm}p{5.0cm}}
\toprule
\textbf{Check} & \textbf{Validation Rule} & \textbf{Interpretation if Failed}\\
\midrule
Missingness & Required fields are non-null (\texttt{session\_id}, event timestamp, metric-specific columns). & Report as undercount risk; do not claim strong support from affected metric.\\
Consistency & Event semantics are stable, for example a cart-add is always a row in \texttt{selected\_items}. & Re-map metric definitions before comparing across sections.\\
Duplicate events & No duplicated session-event records survive deduplication. & Treat inflated rates as invalid until corrected.\\
Sample balance & Cohort is not dominated by a single \texttt{path\_mode} or user subgroup. & Downgrade claim strength to ``partially supported'' or ``inconclusive''.\\
\bottomrule
\end{tabular}
\caption{Validity gates applied before reporting an aggregate metric}
\label{tab:data_validity_checks}
\end{table}

A metric is treated as trustworthy only when it passes the checks in Table~\ref{tab:data_validity_checks}; otherwise the surrounding prose explicitly weakens the claim. Sample-size caveats are repeated in every subsection where the post-filter cohort is small.


\section{Cohort and Integrity}
\label{sec:cohort_integrity}

This section describes the population of sessions on which §5.3--§5.6 build. Every later number inherits the cohort defined here, sourced from \texttt{user\_sessions} together with the per-event tables (\texttt{product\_impressions}, \texttt{user\_ratings}) used to define the active subsets. The numbers reported below are computed at the freeze date 2026-05-09.

\subsection{Cohort Counts}
\label{subsec:cohort_counts}

\begin{table}[H]
\centering
\begin{tabular}{p{6.5cm}p{2.0cm}p{4.5cm}}
\toprule
\textbf{Cohort definition} & \textbf{Sessions} & \textbf{Inclusion filter}\\
\midrule
Raw cohort (open window) & 53 & all rows of \texttt{user\_sessions}\\
Active cohort (had at least one impression) & 25 & inner join with \texttt{product\_impressions}\\
Rated cohort (had at least one rating) & 9 & inner join with \texttt{user\_ratings}\\
Demographic cohort (declared gender) & 50 & \texttt{gender IS NOT NULL}\\
Order-completing cohort (\texttt{ended\_by='order'}) & 8 & \texttt{user\_sessions.ended\_by} column\\
\bottomrule
\end{tabular}
\caption{Cohort counts at freeze date 2026-05-09 (open window, earliest session 2026-04-20)}
\label{tab:cohort_counts}
\end{table}

The active cohort (25 of 53, 47.2\%) is the most stable subset for behavioural metrics: it represents sessions in which the user actually saw search results. The rated cohort (9 of 53, 17.0\%) is small in absolute terms; §5.6 therefore reports its findings as anecdotal evidence and annotates every aggregate with the sample size $N=9$.

\subsection{Path-Mode Sample Balance}
\label{subsec:path_balance}

\begin{table}[H]
\centering
\begin{tabular}{p{3cm}p{3cm}p{6cm}}
\toprule
\textbf{path\_mode} & \textbf{Sessions} & \textbf{Notes}\\
\midrule
\texttt{path1} & 35 & Primary text-to-image search path; the agent reasoning loop runs for these sessions\\
\texttt{path2} & 7  & Image upload path (single-stage cosine lookup, no agent reasoning); secondary in this chapter\\
\bottomrule
\end{tabular}
\caption{Sessions appearing in any per-event table partitioned by \texttt{path\_mode}}
\label{tab:path_balance}
\end{table}

Some sessions span both paths (the user uploaded an image and also typed a query). PATH~1 dominates the active cohort; PATH~2 is reported as a brief secondary table in §5.5 because its lower volume does not support standalone claims about agent quality.

\subsection{Funnel Integrity Flags}
\label{subsec:funnel_integrity}

The function \texttt{\_evaluate\_funnel\_integrity} in \texttt{agent/memory.py} flags four logical inconsistencies per session: clicks without prior impressions, cart-adds without prior impressions, intent events without a prior cart-add, and orders without a prior cart-add. Aggregating across the active cohort, no sessions in the present window failed any of these checks; the funnel data is therefore internally consistent and the §5.5 aggregates do not need to be downgraded for integrity reasons.



\section{Knowledge-Base Quality}
\label{sec:kb_quality}

This section reports on the catalogue and its enrichment, sourced from \texttt{fashion\_items} joined with \texttt{fashion\_item\_enrichment}. The relevant claim from Chapter~3 (Section~3.2) is that the multimodal enrichment pass produces a usable knowledge base for downstream vector search. The evidence below supports that claim by demonstrating complete coverage and caption length within the targeted band.

\subsection{Catalogue Size and Enrichment Coverage}
\label{subsec:enrichment_coverage}

\begin{table}[H]
\centering
\begin{tabular}{p{6.5cm}p{3cm}p{3cm}}
\toprule
\textbf{Quantity} & \textbf{Count} & \textbf{Coverage}\\
\midrule
Total items in \texttt{fashion\_items} & 5,096 & ---\\
Enrichment rows in \texttt{fashion\_item\_enrichment} & 5,096 & 100.0\%\\
Items with non-empty \texttt{caption} & 5,096 & 100.0\%\\
Items with non-empty \texttt{color} & 5,096 & 100.0\%\\
\bottomrule
\end{tabular}
\caption{Catalogue size and enrichment coverage at freeze date 2026-05-09}
\label{tab:enrichment_coverage}
\end{table}

The enrichment pass has covered the catalogue in full. There is no missing-caption or missing-color tail, so the downstream indexing pass (§3.2) had a complete input.

\subsection{Caption Length Distribution}
\label{subsec:caption_length}

The Gemini caption prompt targets 30--40 words per caption. Across the 5,096 enriched items, the empirical distribution of caption length, measured in whitespace-separated tokens, has the following statistics:

\begin{table}[H]
\centering
\begin{tabular}{p{4.5cm}p{3cm}}
\toprule
\textbf{Statistic} & \textbf{Value (words)}\\
\midrule
Minimum & 17\\
Mean & 34.9\\
Maximum & 48\\
Target band stated in §3.2 & 30--40\\
\bottomrule
\end{tabular}
\caption{Caption-length statistics over 5,096 enriched items}
\label{tab:caption_length}
\end{table}

The mean of 34.9 words sits inside the 30--40 target band. The minimum of 17 words and the maximum of 48 words indicate that the enrichment pass is not strictly bounded, but no caption falls below half the lower target nor exceeds the upper target by more than 20\%. In practical terms, the enrichment output is at the intended length for the search captions ingested by the SigLIP text encoder.

\subsection{Catalogue Category Distribution}
\label{subsec:category_distribution}

\begin{table}[H]
\centering
\begin{tabular}{p{3cm}p{2cm}p{3cm}p{2cm}}
\toprule
\textbf{Category} & \textbf{Count} & \textbf{Category} & \textbf{Count}\\
\midrule
T-Shirt    & 1011 & Shorts     & 308\\
Longsleeve & 699  & Hat        & 171\\
Pants      & 692  & Skirt      & 155\\
Shoes      & 431  & Polo       & 120\\
Shirt      & 378  & Undershirt & 118\\
Dress      & 357  & Blazer     & 109\\
Outwear    & 312  & Hoodie     & 100\\
\bottomrule
\end{tabular}
\caption{Top fourteen categories by item count in \texttt{fashion\_items}}
\label{tab:category_distribution}
\end{table}

The catalogue is dominated by upper-body garments (T-Shirt, Longsleeve, Shirt, Polo) and bottoms (Pants, Shorts, Skirt). The retrieval pipeline in §3.4 is therefore exercised mostly on those categories; categories with under 100 items are present but contribute marginally to retrieval traffic.




\section{LLM Call Discipline}
\label{sec:llm_call_discipline}

This section reports the distribution of LLM calls per turn, sourced from \texttt{llm\_token\_usage} together with the views \texttt{session\_token\_summary} and \texttt{mode\_cost\_summary}. The corresponding claim in Chapter~3 (Section~3.5) is that the deployed pipeline issues at most two LLM calls per turn (intent classification, then synthesis), and never enters a ReAct-style loop. The evidence below confirms that empirically.

\subsection{Call-Name Distribution}
\label{subsec:call_name_distribution}

\begin{table}[H]
\centering
\begin{tabular}{p{3cm}p{2.5cm}p{3.5cm}p{3.5cm}}
\toprule
\textbf{call\_name} & \textbf{Rows} & \textbf{Avg. input tokens} & \textbf{Avg. output tokens}\\
\midrule
\texttt{intent}    & 155 & 1{,}998 & 120\\
\texttt{synthesis} & 60  & 859     & 223\\
\texttt{expansion} & 0   & ---     & ---\\
\bottomrule
\end{tabular}
\caption{Per-call token usage across the cohort, from \texttt{llm\_token\_usage}}
\label{tab:call_name_distribution}
\end{table}

Three observations follow from Table~\ref{tab:call_name_distribution}:

\begin{enumerate}
    \item \textbf{Every turn invokes the intent classifier.} The 155 \texttt{intent} rows correspond to 155 turns, and they are the only mandatory call.
    \item \textbf{Synthesis is conditional, not mandatory.} Only 60 of the 155 turns proceeded to a synthesis call. The ratio $60/155 \approx 38.7\%$ is the fraction of user turns that produced product results worth synthesising; the remaining 61.3\% were either clarification responses (no synthesis call needed, the agent emits a templated question) or out-of-scope requests (refused without retrieval).
    \item \textbf{The query-expansion call is never recorded.} The short-query gate at six words (Algorithm~\ref{alg:expand_query}) suppressed expansion in every observed turn; \texttt{call\_name='expansion'} has zero rows. In practice the deployed cohort always reached search with the user's original query.
\end{enumerate}

The empirical per-turn LLM-call count is therefore exactly one (intent only, for clarification or out-of-scope turns) or exactly two (intent then synthesis, for product turns). No turn invoked three or more LLM calls. This rules out, by direct counting, any ReAct-style multi-call loop in the deployed cohort.

\subsection{Per-Session Token Aggregates}
\label{subsec:token_aggregates}

\begin{table}[H]
\centering
\begin{tabular}{p{6cm}p{4cm}}
\toprule
\textbf{Quantity} & \textbf{Value}\\
\midrule
Sessions with at least one logged call & 41\\
Average input tokens per session & 8{,}811\\
Average output tokens per session & 778\\
Average total tokens per session & 9{,}590\\
\bottomrule
\end{tabular}
\caption{Per-session token totals from \texttt{session\_token\_summary}}
\label{tab:token_aggregates}
\end{table}

\subsection{Cost per Turn}
\label{subsec:cost_per_turn}

\begin{table}[H]
\centering
\begin{tabular}{p{3cm}p{3cm}p{2cm}p{2cm}p{3cm}}
\toprule
\textbf{Orchestration} & \textbf{Synthesizer} & \textbf{Sessions} & \textbf{Turns} & \textbf{Avg. USD/turn}\\
\midrule
\texttt{direct} & \texttt{gemini-2.5-flash} & 24 & 50 & \$0.000132\\
\texttt{direct} & \texttt{gemini-2.5-flash} (legacy ``fixed'' label) & 6 & 10 & \$0.000129\\
\bottomrule
\end{tabular}
\caption{Estimated per-turn USD cost from \texttt{mode\_cost\_summary}}
\label{tab:cost_per_turn}
\end{table}

A representative turn costs approximately \$0.00013, that is, on the order of \$0.13 per one thousand turns at the published Gemini~2.5~Flash list price. The cost is dominated by the intent-classifier input (large system prompt) rather than by the synthesis output. The two rows differ only in a legacy \texttt{orchestrator\_model} label (\texttt{fixed} versus \texttt{gemini-2.5-flash}) emitted by an earlier version of the agent; both rows correspond to the same direct-routing pipeline.


\section{Implicit Relevance from User Behaviour}
\label{sec:implicit_relevance}

This section uses behavioural events to argue that the items returned by the retrieval pipeline are perceived as relevant by users. The argument is implicit, not gold-set-based: if the retrieval pipeline returned nonsense, users would not click on the results, would not add them to the cart, and would not signal will-buy. Observed click, cart-add, and will-buy rates therefore act as a behavioural floor on retrieval relevance, sourced from \texttt{product\_impressions}, \texttt{product\_clicks}, \texttt{selected\_items}, \texttt{cart\_removals}, \texttt{product\_intents}, and the view \texttt{session\_path\_funnel\_summary}.

\subsection{PATH 1 Funnel}
\label{subsec:path1_funnel}

\begin{table}[H]
\centering
\begin{tabular}{p{3.5cm}p{2cm}p{4.5cm}}
\toprule
\textbf{Step} & \textbf{Count} & \textbf{Conversion rate to next step}\\
\midrule
Sessions in funnel  & 30 & ---\\
Impressions logged   & 254 & 8.5 impressions per funnel session (mean)\\
Clicks               & 81 & 31.9\% of impressions (CTR)\\
Cart-adds            & 21 & 25.9\% of clicks; 8.3\% of impressions\\
Will-buy intents     & 2 & 9.5\% of cart-adds\\
Not-for-me intents   & 1 & 4.8\% of cart-adds\\
Orders (session-end) & 9 & ---\\
\bottomrule
\end{tabular}
\caption{PATH 1 behavioural funnel from \texttt{session\_path\_funnel\_summary}}
\label{tab:path1_funnel}
\end{table}

The click-through rate of 31.9\% is the strongest signal in the funnel: roughly one in three displayed items receives an explicit user click. The cart-add rate (8.3\% of impressions) drops by a factor of four from the click rate, which is consistent with users browsing freely and committing more selectively. The will-buy intent rate is small in absolute terms (only two events from 21 cart-adds), reflecting both genuine selectivity and the small absolute cohort; it is therefore reported but not used as a primary quality signal.

\subsection{PATH 2 Funnel (secondary)}
\label{subsec:path2_funnel}

\begin{table}[H]
\centering
\begin{tabular}{p{3.5cm}p{2cm}p{4.5cm}}
\toprule
\textbf{Step} & \textbf{Count} & \textbf{Conversion rate}\\
\midrule
Sessions in funnel & 7 & ---\\
Impressions logged & 42 & ---\\
Clicks             & 6 & 14.3\% of impressions\\
Cart-adds          & 2 & 33.3\% of clicks; 4.8\% of impressions\\
Will-buy intents   & 0 & ---\\
Orders             & 0 & ---\\
\bottomrule
\end{tabular}
\caption{PATH 2 behavioural funnel; included for completeness, not used to support claims}
\label{tab:path2_funnel}
\end{table}

PATH~2's small cohort (7 sessions) and complete absence of will-buy or order events mean the table is reported here for completeness only. PATH~2's funnel does not support any claim about agent quality, as discussed in §5.1 sample-balance.

\subsection{Position of Cart-Adds (Reranker Evidence)}
\label{subsec:position_histogram}

If the BGE cross-encoder rerank in §3.4 places relevant items near the top of the result list, cart-add events should concentrate at low \texttt{position} values. The distribution of \texttt{selected\_items.position} in the PATH 1 cohort is:

\begin{table}[H]
\centering
\begin{tabular}{p{2cm}p{2cm}p{8cm}}
\toprule
\textbf{Position} & \textbf{Count} & \textbf{Notes}\\
\midrule
0 & 4  & 19.05\% --- cart-add via numeric selection without prior click\\
1 & 5  & top-of-list\\
2 & 3 & \\
3 & 2 & \\
4 & 2 & \\
5 & 3 & \\
6 & 2 & last visible position\\
\bottomrule
\end{tabular}
\caption{PATH 1 cart-add position distribution from \texttt{selected\_items.position}}
\label{tab:position_histogram}
\end{table}

The share of \texttt{position=0} is 19.05\%, just below the 20\% threshold beyond which the histogram would have been recomputed from the impressions join (the methodological design decision recorded in the change proposal). The distribution among positions 1--6 is tightly clustered with five cart-adds at position~1 (the modal position) and a gentle tail. This is consistent with the reranker putting useful items near the top, but the absolute counts (21 cart-adds in total) are too small to be claimed as statistically significant; the histogram is reported as suggestive evidence, not as a confirmatory test.


\section{Subjective Quality}
\label{sec:subjective_quality}

This section reports on the rating signals collected from \texttt{user\_ratings}. Because $N = 9$, every aggregate below is presented as anecdotal evidence and not as a statistical claim. The intent of this section is to triangulate the implicit signals in §5.5 with explicit user feedback.

\subsection{Rating Distributions}
\label{subsec:rating_distributions}

\begin{table}[H]
\centering
\begin{tabular}{p{4cm}p{2cm}p{2cm}p{4cm}}
\toprule
\textbf{Rating dimension} & \textbf{Mean} & \textbf{SD} & \textbf{N}\\
\midrule
\texttt{rating\_overall}        & 4.56 & 0.53 & 9\\
\texttt{rating\_suggestions}    & 4.67 & 0.50 & 9\\
\texttt{rating\_conversation}   & 4.22 & 0.83 & 9\\
\bottomrule
\end{tabular}
\caption{User ratings on a 1--5 scale, $N=9$, freeze date 2026-05-09}
\label{tab:rating_distributions}
\end{table}

\begin{table}[H]
\centering
\begin{tabular}{p{3cm}p{2cm}p{6cm}}
\toprule
\textbf{rating\_overall} & \textbf{Count} & \textbf{Notes}\\
\midrule
1 & 0 & ---\\
2 & 0 & ---\\
3 & 0 & ---\\
4 & 4 & \\
5 & 5 & \\
\bottomrule
\end{tabular}
\caption{Distribution of \texttt{rating\_overall} among the rated cohort}
\label{tab:rating_overall_distribution}
\end{table}

Within the rated cohort, every score for \texttt{rating\_overall} is either~4 or~5; there are no scores of~1, 2, or~3. The means for the three dimensions sit between 4.2 and 4.7, with the conversation rating slightly lower than the suggestions and overall ratings. The sample is small and self-selected: only users who chose to leave a rating appear in the table, and that subset may be biased toward more engaged or more positive users. The numbers are therefore presented as descriptive --- they are consistent with the implicit relevance signals in §5.5, but they do not establish a causal claim about agent quality.

\subsection{Free-Text Feedback Coverage}
\label{subsec:feedback_coverage}

Of the 9 ratings, 8 (88.9\%) include a non-empty \texttt{feedback} string. Of the 3 \texttt{product\_intents} rows in the cohort, 0 (0.0\%) include a non-empty \texttt{reason} string. A manual coding of feedback themes is therefore feasible at the rating level (with $N=8$ free-text samples) but not at the per-item intent level (no reasons recorded). Because the rating sample is so small, no thematic table is produced; representative feedback excerpts are not transcribed here to preserve respondent anonymity.

\subsection{Limits of the Subjective Evidence}
\label{subsec:subjective_limits}

Three caveats apply to every number in §5.6:

\begin{enumerate}
    \item Sample size $N=9$ falls below conventional thresholds for statistical inference; no test of mean differences is performed.
    \item Self-selection bias: users who rated may differ systematically from users who did not; the rated cohort is 17.0\% of the raw cohort.
    \item No paired design: there is no matched pre-deployment baseline against which to compare these scores. A mean of~4.56 on a 1--5 scale is favourable in absolute terms but cannot be argued as ``better than the previous system'' without comparable historical data.
\end{enumerate}

These caveats motivate the position taken in §5.6.1 of presenting the ratings as descriptive triangulation rather than as primary evidence.
